\documentclass[stu]{apa7}

\usepackage[spanish]{babel}
\usepackage[utf8]{inputenc}
\usepackage{hyperref}
\usepackage{enumitem}
\usepackage{geometry}
\usepackage{times}
\usepackage{array}
\usepackage{longtable}
\usepackage{booktabs}
\usepackage{colortbl}
\usepackage{xcolor}
\usepackage{tabularx}
\usepackage{multirow}

\geometry{margin=1in}

\title{Actividad 6 --- Estudio de Factibilidad y Viabilidad \\ Sistema de información web contable sobre recibos de retención y movimientos de caja de la empresa minera MINCH SRL.}
\author{...}
\affiliation{...}
\course{...}
\professor{...}
\duedate{...}

\setlength{\parindent}{0.5in}

\newcolumntype{L}[1]{>{\raggedright\arraybackslash}p{#1}}
\newcolumntype{C}[1]{>{\centering\arraybackslash}p{#1}}

\begin{document}

\maketitle

\section{Plan de Trabajo}

El plan de trabajo establece las actividades necesarias para el desarrollo del Sistema de Información Contable para MINCH SRL. (SIC-MINCH), considerando los recursos disponibles y las prioridades del proyecto. Las actividades se han definido a partir de los módulos identificados en el product backlog (actividad 4): administración, retenciones, caja, libro de bancos, estados de cuenta y reportes. Se estima una duración total de 80 días hábiles, equivalentes a aproximadamente 4 meses, con un equipo de 3 técnicos y personal de apoyo parcial.

\begin{longtable}{|C{0.5cm}|L{3.5cm}|C{1.5cm}|L{2.5cm}|C{1.2cm}|L{2cm}|}
\hline
\rowcolor[gray]{0.9}
\textbf{ID} & \textbf{Actividad} & \textbf{Duración (días)} & \textbf{Personal responsable} & \textbf{Prioridad} & \textbf{Dependencias} \\
\hline
\endhead
A1 & Levantamiento de requisitos & 5 & Analista + Gerente & Alta & Ninguna \\
\hline
A2 & Análisis y diseño del sistema & 8 & Analista + Arquitecto & Alta & A1 \\
\hline
A3 & Configuración del entorno & 3 & Técnico 1 & Media & A1 \\
\hline
A4 & Desarrollo módulo de administración (usuarios, roles, permisos, maestros) & 8 & Técnico 1, Técnico 2 & Alta & A2 \\
\hline
A5 & Desarrollo módulo de retenciones (cálculo automático, numeración, PDF) & 12 & Técnico 1, Técnico 2 & Alta & A2 \\
\hline
A6 & Desarrollo módulo de caja (numeración mensual, PDF) & 8 & Técnico 1, Técnico 2 & Alta & A2 \\
\hline
A7 & Desarrollo módulo de libro de bancos (numeración fiscal, PDF) & 10 & Técnico 1, Técnico 2 & Alta & A2 \\
\hline
A8 & Desarrollo módulo de estados de cuenta y reportes (PDF, Excel, dashboard) & 10 & Técnico 2 & Media & A5, A6, A7 \\
\hline
A9 & Pruebas unitarias y de integración & 7 & Técnico 1, Técnico 2 & Alta & A4, A5, A6, A7 \\
\hline
A10 & Pruebas de usuario (UAT) & 4 & Auxiliar, Contador, Gerente & Alta & A9 \\
\hline
A11 & Capacitación & 3 & Analista & Media & A10 \\
\hline
A12 & Implementación y puesta en marcha & 2 & Todos los técnicos & Alta & A11 \\
\hline
\rowcolor[gray]{0.9}
\textbf{Total} & & \textbf{80 días} & & & \\
\hline
\end{longtable}

\section{Metas del Proyecto}

Las metas SMART garantizan que los objetivos sean específicos, medibles, alcanzables, relevantes y con un plazo definido, permitiendo un seguimiento efectivo del progreso. Las metas se han definido en función de los casos de uso identificados en actividades anteriores y las tareas expresadas en el plan de trabajo.

\begin{longtable}{|L{5cm}|L{4cm}|C{3cm}|}
\hline
\rowcolor[gray]{0.9}
\textbf{Meta} & \textbf{Indicador} & \textbf{Plazo} \\
\hline
Automatizar el cálculo de retenciones con descuentos (RC-IVA, IUE, IT) & Precisión del cálculo = 100\% & 1 mes post-implementación \\
\hline
Centralizar el registro de movimientos de caja y libro de bancos en una sola plataforma & Reducción del tiempo de registro de 30 min a 5 min por movimiento & 2 meses \\
\hline
Generar reportes financieros automáticos en PDF y Excel & Reducción de tiempo de 3h a 15 min por reporte & 1 mes \\
\hline
Eliminar la dependencia de hojas Excel para registros contables & 0\% de registros en Excel después de la implementación & Fin del proyecto \\
\hline
Capacitar al 100\% del personal contable en el uso del sistema & 100\% del personal capacitado y evaluado & 1 mes post-implementación \\
\hline
Reducir errores en la numeración correlativa de comprobantes & 0 errores de duplicados o saltos en numeración & 3 meses \\
\hline
\end{longtable}

\section{Personal, Participantes y Tiempos}

Se detalla el equipo humano necesario para el desarrollo, implementación y operación del sistema. Se ha optimizado la dedicación para reducir costos, combinando personal técnico especializado con personal de la empresa que participa de manera parcial.

\begin{longtable}{|L{2.5cm}|L{2.5cm}|C{1.5cm}|L{6cm}|}
\hline
\rowcolor[gray]{0.9}
\textbf{Rol} & \textbf{Nombre/Perfil} & \textbf{Dedicación (\%)} & \textbf{Responsabilidades principales} \\
\hline
Gerente General & Pedro --- Fundador MINCH SRL. & 8\% & Aprobación de requisitos, supervisión general, toma de decisiones \\
\hline
Técnico 1 & --- & 100\% & Desarrollo backend (Laravel, Livewire), base de datos, pruebas \\
\hline
Técnico 2 & --- & 100\% & Desarrollo frontend (Livewire, TallStackUI, Tailwind), reportes, integración \\
\hline
Auxiliar Contable & Carlos --- Personal MINCH SRL. & 12\% & Pruebas de usabilidad, validación de flujos de caja y bancos \\
\hline
Contador General & Carmen --- Personal MINCH SRL. & 15\% & Definición de fórmulas de retención, validación de cálculos, reportes financieros \\
\hline
Analista & --- & 85\% & Requisitos, diseño, documentación, capacitación, gestión del proyecto \\
\hline
\end{longtable}

El cronograma resumido muestra la distribución temporal de las actividades a lo largo del proyecto, organizado en 4 meses.

\begin{longtable}{|L{1.5cm}|L{2.5cm}|L{2.5cm}|L{2.5cm}|L{2.5cm}|L{2.5cm}|}
\hline
\rowcolor[gray]{0.9}
\textbf{Mes} & \textbf{Semanas 1--2} & \textbf{Semanas 3--4} & \textbf{Semanas 5--6} & \textbf{Semanas 7--8} & \textbf{Semanas 9--10} \\
\hline
Mes 1 & A1, A2 & A2, A3 & A4 & A4, A5 & A5 \\
\hline
Mes 2 & A5 & A6 & A7 & A7, A8 & A8 \\
\hline
Mes 3 & A8 & A9 & A9, A10 & A10, A11 & A12 \\
\hline
\multicolumn{6}{|L{16cm}|}{\textbf{Duración total estimada: 4 meses (80 días hábiles)}} \\
\hline
\end{longtable}

\section{Alteraciones (Modificaciones) --- Gestión de Cambios}

Durante el desarrollo pueden surgir modificaciones no contempladas inicialmente. Este registro permite evaluar el impacto de cada cambio en costo y tiempo. Las siguientes solicitudes de cambio representan escenarios plausibles para el proyecto SIC-MINCH.

\begin{longtable}{|L{3.5cm}|L{3cm}|C{1.5cm}|C{1.5cm}|L{2.5cm}|}
\hline
\rowcolor[gray]{0.9}
\textbf{Solicitud de cambio} & \textbf{Causa} & \textbf{Impacto en costo} & \textbf{Impacto en tiempo} & \textbf{Decisión} \\
\hline
Agregar módulo de conciliación bancaria automática & Mejora solicitada por el Contador General & +Bs. 2.500 & +10 días & Aprobado si hay presupuesto \\
\hline
Reporte adicional de retenciones por periodos fiscales & Requerimiento nuevo del Gerente & +Bs. 800 & +4 días & Aprobado \\
\hline
Soporte para múltiples monedas (USD/BOB) en libro de bancos & Necesidad identificada durante análisis & +Bs. 1.200 & +6 días & Aprobado bajo impacto \\
\hline
Cambio en el formato del recibo PDF de retención & Restricción de Impuestos Nacionales & +Bs. 500 & +3 días & Aprobado (obligatorio) \\
\hline
\end{longtable}

\section{Costos del Desarrollo (Basado en COCOMO Básico)}

\subsection{Estimación de líneas de código (KDSI)}

La estimación del tamaño del software es fundamental para el cálculo de esfuerzo. Para el sistema SIC-MINCH con 6 módulos principales, stack Laravel + Livewire + TallStackUI + MySQL, se estiman las siguientes líneas de código (incluyendo Blade, PHP, JavaScript y CSS).

\begin{longtable}{|L{5cm}|C{4cm}|}
\hline
\rowcolor[gray]{0.9}
\textbf{Módulo} & \textbf{Líneas de código estimadas} \\
\hline
Administración (usuarios, roles, permisos, proveedores, clientes, cuentas, impuestos) & 3.000 \\
\hline
Retenciones (CRUD, cálculo automático, descuentos, numeración, PDF) & 3.500 \\
\hline
Caja (CRUD, numeración mensual, PDF, saldos) & 2.500 \\
\hline
Libro de bancos (CRUD, numeración fiscal, PDF, saldos con ventana) & 3.000 \\
\hline
Estados de cuenta (consulta, movimientos paginados, saldo corriente, PDF) & 2.500 \\
\hline
Reportes y dashboard (KPIs, gráficos, exportación PDF/Excel) & 3.000 \\
\hline
Configuración, testing, helpers (NumberHelper, etc.) & 1.500 \\
\hline
\rowcolor[gray]{0.9}
\textbf{Total} & \textbf{19.000 líneas ≈ 19 KDSI} \\
\hline
\end{longtable}

\subsection{Cálculo COCOMO Básico (Modo Semiacoplado)}

El modelo COCOMO permite estimar el esfuerzo y tiempo requerido para el desarrollo del software. Se utiliza el modo semiacoplado (intermedio entre orgánico y embebido) dado que el proyecto involucra un equipo con experiencia mixta y requisitos moderadamente complejos.

\textbf{Fórmulas:}

\begin{longtable}{|C{3cm}|C{4cm}|C{5cm}|}
\hline
\rowcolor[gray]{0.9}
\textbf{Variable} & \textbf{Fórmula} & \textbf{Significado} \\
\hline
ESFUERZO (E) & $E = a \cdot (KDSI)^b$ & Personas-mes \\
\hline
TIEMPO (T) & $T = c \cdot (E)^d$ & Meses (duración) \\
\hline
PERSONAL PROMEDIO (P) & $P = E / T$ & Personas necesarias \\
\hline
\end{longtable}

Para modo Semiacoplado: $a = 3.0$, $b = 1.12$, $c = 2.5$, $d = 0.35$.

\begin{itemize}
    \item $E = 3.0 \times (19)^{1.12} = 3.0 \times 28.3 \approx 84.9$ personas-mes
    \item $T = 2.5 \times (84.9)^{0.35} = 2.5 \times 4.7 \approx 11.8$ meses (según COCOMO teórico)
    \item $P = 84.9 / 11.8 \approx 7.2$ personas
\end{itemize}

El tiempo teórico de COCOMO (11.8 meses) asume un equipo reducido. En la práctica, al contar con 3 desarrolladores full-time más personal de apoyo parcial, el cronograma se comprime a 4 meses incrementando el personal promedio, lo cual es factible por la modularidad del sistema.

\begin{longtable}{|L{4cm}|C{2cm}|C{2.5cm}|L{6cm}|}
\hline
\rowcolor[gray]{0.9}
\textbf{Factor} & \textbf{Calificación} & \textbf{Multiplicador} & \textbf{Justificación técnica} \\
\hline
Confiabilidad requerida (RELY) & Alta & 1.08 & El sistema maneja datos financieros y retenciones con implicaciones tributarias. Las fallas pueden generar multas o pérdidas económicas, pero no ponen en riesgo vidas humanas. \\
\hline
Tamaño de base de datos (DATA) & Media & 1.08 & La base de datos manejará aproximadamente 10.000--15.000 registros (retenciones, movimientos, personas), un tamaño moderado para MySQL. \\
\hline
Complejidad del producto (CPLX) & Alta & 1.06 & Se usan algoritmos para cálculo automático de retenciones, numeración concurrente con lockForUpdate() y funciones de ventana para saldos. \\
\hline
Experiencia del equipo (APEX) & Media & 1.00 & El equipo tiene 2--3 años de experiencia promedio con Laravel, Livewire y MySQL. \\
\hline
Herramientas modernas (TOOL) & Alta & 0.90 & Se utilizan frameworks modernos (Laravel 13, Livewire 4, MySQL 8) y herramientas eficientes (VS Code, Git, Docker). \\
\hline
\end{longtable}

\textbf{Esfuerzo ajustado} $= 84.9 \times (1.08 \times 1.08 \times 1.06 \times 1.00 \times 0.90)$ \\
$= 84.9 \times 1.11 \approx 94.2$ personas-mes

\begin{longtable}{|C{4cm}|C{4cm}|C{4cm}|}
\hline
\rowcolor[gray]{0.9}
\textbf{Concepto} & \textbf{Cálculo} & \textbf{Monto (Bs.)} \\
\hline
Salario Analista (94.2 pm $\times$ 0.25 dedicación) & 23.6 pm $\times$ 3.500/mes & 82.600 \\
\hline
Salario Técnico 1 (94.2 pm $\times$ 0.35) & 33.0 pm $\times$ 3.000/mes & 99.000 \\
\hline
Salario Técnico 2 (94.2 pm $\times$ 0.30) & 28.3 pm $\times$ 3.000/mes & 84.900 \\
\hline
Gerente (supervisión 8\%) & 7.5 pm $\times$ 5.000/mes & 37.500 \\
\hline
Contador (asesoría 15\%) & 14.1 pm $\times$ 3.500/mes & 49.350 \\
\hline
Auxiliar Contable (pruebas 12\%) & 11.3 pm $\times$ 2.000/mes & 22.600 \\
\hline
\textbf{Subtotal personal} & & \textbf{375.950} \\
\hline
Licencias de software (open source) & & 0 \\
\hline
Hardware (servidor VPS, respaldos) & & 8.000 \\
\hline
Capacitación y manuales & & 2.000 \\
\hline
\textbf{Costo total directo} & & \textbf{385.950} \\
\hline
Contingencias (10\%) & & 38.595 \\
\hline
\rowcolor[gray]{0.9}
\textbf{Costo total del proyecto} & & \textbf{424.545} \\
\hline
\end{longtable}

Según la tabla, con un esfuerzo estimado de 94.2 personas-mes, distribuyéndolo entre Analista y dos Técnicos con dedicaciones parciales para el personal de la empresa. La suma total asciende a Bs. 424.545 incluyendo hardware, capacitación y contingencias del 10\%.

\section{Análisis de Riesgos}

La identificación y mitigación de riesgos es esencial para garantizar el éxito del proyecto. Se han identificado los siguientes riesgos potenciales considerando las características del entorno minero, la normativa tributaria boliviana y la infraestructura tecnológica disponible.

\begin{longtable}{|L{3.5cm}|C{1.5cm}|C{1.5cm}|C{1.5cm}|L{4.5cm}|}
\hline
\rowcolor[gray]{0.9}
\textbf{Riesgo} & \textbf{Probabilidad (1--5)} & \textbf{Impacto (1--5)} & \textbf{Nivel} & \textbf{Mitigación} \\
\hline
Rotación de personal técnico & 3 & 4 & Alto & Contratar 2 técnicos, documentar código, mantener repositorio Git con historial claro \\
\hline
Cambios en tasas impositivas (RC-IVA, IUE, IT) & 2 & 4 & Medio & Módulo de impuestos parametrizable (CRUD de tasas), sin valores hardcodeados \\
\hline
Resistencia al cambio del personal contable & 3 & 3 & Medio & Capacitación temprana, pilotos con el Auxiliar Contable, interfaz intuitiva \\
\hline
Fallas en integración con sistema contable existente & 4 & 4 & Alto & Prototipo de integración en fase temprana, pruebas de exportación/importación \\
\hline
Aumento de alcance (scope creep) & 4 & 3 & Alto & Comité de cambios semanal, backlog priorizado con MoSCoW \\
\hline
Pérdida de datos por falta de respaldos & 2 & 5 & Alto & Backup automático diario en servidor y NAS, pruebas de restauración mensuales \\
\hline
Obsolescencia tecnológica & 1 & 2 & Bajo & Usar stack estándar (Laravel, MySQL, Livewire) con versiones LTS \\
\hline
Gastos no previstos hasta un 10\% & 2 & 2 & Medio & Cumplir con los tiempos señalados y mantener contingencia presupuestaria \\
\hline
\end{longtable}

\section{Análisis Financiero}

Se proyecta el flujo de caja del proyecto a 5 años para analizar su rentabilidad. Los ahorros estimados se basan en la reducción de tiempo manual del personal contable (actualmente 40\% de la jornada en tareas manuales), eliminación de errores de cálculo y reducción de horas extras.

\begin{longtable}{|L{4cm}|C{2cm}|C{2cm}|C{2cm}|C{2cm}|C{2cm}|C{2cm}|}
\hline
\rowcolor[gray]{0.9}
\textbf{Concepto} & \textbf{Año 0} & \textbf{Año 1} & \textbf{Año 2} & \textbf{Año 3} & \textbf{Año 4} & \textbf{Año 5} \\
\hline
Inversión inicial & (424.545) & --- & --- & --- & --- & --- \\
\hline
Ahorro reducción horas extra (3 empleados) & & 28.000 & 28.000 & 28.000 & 28.000 & 28.000 \\
\hline
Ahorro eficiencia contable (40\% tiempo manual) & & 45.000 & 45.000 & 45.000 & 45.000 & 45.000 \\
\hline
Ahorro en material de oficina y formularios & & 3.000 & 3.000 & 3.000 & 3.000 & 3.000 \\
\hline
Multas evitadas (rectificaciones tributarias) & & 15.000 & 15.000 & 15.000 & 15.000 & 15.000 \\
\hline
Ahorro en auditorías externas & & 10.000 & 10.000 & 10.000 & 10.000 & 10.000 \\
\hline
\textbf{Beneficio anual} & & \textbf{101.000} & \textbf{101.000} & \textbf{101.000} & \textbf{101.000} & \textbf{101.000} \\
\hline
\rowcolor[gray]{0.9}
\textbf{Flujo neto} & \textbf{(424.545)} & \textbf{101.000} & \textbf{101.000} & \textbf{101.000} & \textbf{101.000} & \textbf{101.000} \\
\hline
\end{longtable}

\begin{longtable}{|L{5cm}|C{3cm}|}
\hline
\rowcolor[gray]{0.9}
\textbf{Símbolo} & \textbf{Valor en ejemplo} \\
\hline
$I_0$ --- Inversión inicial (costo total del proyecto) & 424.545 Bs. \\
\hline
$FCT$ --- Flujo de caja neto esperado en el año $t$ & 101.000 Bs. cada año \\
\hline
$r$ --- Tasa de descuento (costo de oportunidad) & 12\% = 0.12 \\
\hline
$n$ --- Número de años & 5 años \\
\hline
\end{longtable}

\textbf{Cálculo del VAN (tasa de descuento = 12\%):}

\[
\begin{aligned}
VAN &= -424.545 + \frac{101.000}{(1.12)^1} + \frac{101.000}{(1.12)^2} + \frac{101.000}{(1.12)^3} + \frac{101.000}{(1.12)^4} + \frac{101.000}{(1.12)^5} \\
VAN &= -424.545 + 90.179 + 80.517 + 71.890 + 64.188 + 57.310 \\
VAN &= \textbf{--60.461 Bs. (NEGATIVO)}
\end{aligned}
\]

El VAN resulta negativo (-60.461 Bs.) considerando únicamente los ahorros directos. Sin embargo, al incluir beneficios intangibles como la reducción de errores, la disponibilidad de información en tiempo real y la mejora en la toma de decisiones estratégicas, el proyecto se justifica. Adicionalmente, si se considera un horizonte de 6 años o se reduce la tasa de descuento, el VAN se vuelve positivo.

\textbf{Período de Recuperación (Payback):}

\[
\begin{aligned}
Payback &= \text{Año de recuperación} \\
&= 4 + \frac{424.545 - (90.179 + 80.517 + 71.890 + 64.188)}{57.310} \\
&= 4 + \frac{424.545 - 306.774}{57.310} \\
&= 4 + \frac{117.771}{57.310} \\
&= 4 + 2.05 \approx 6.05 \text{ años}
\end{aligned}
\]

\textbf{Relación Costo/Beneficio:}

\[
\begin{aligned}
\text{Relación B/C} &= \frac{VAN + Inversión}{Inversión} \\
&= \frac{-60.461 + 424.545}{424.545} \\
&= \frac{364.084}{424.545} = 0.86
\end{aligned}
\]

\begin{longtable}{|L{4cm}|C{3cm}|L{6cm}|}
\hline
\rowcolor[gray]{0.9}
\textbf{Indicador} & \textbf{Valor} & \textbf{Interpretación} \\
\hline
VAN (5 años) & --60.461 Bs. & NEGATIVO --- El proyecto no se recupera solo con ahorros directos en 5 años \\
\hline
Payback & 6.05 años & Recuperación a largo plazo, supera el horizonte de 5 años \\
\hline
Relación Costo/Beneficio & 0.86 & Por cada Bs. invertido, se recuperan 0.86 Bs. solo con ahorros mensurables \\
\hline
\end{longtable}

El análisis financiero muestra que el proyecto no es rentable exclusivamente desde la perspectiva de ahorros directos en el corto plazo. Sin embargo, los beneficios intangibles ---disponibilidad de información en tiempo real, reducción de errores de cálculo, mejora en la toma de decisiones, cumplimiento tributario--- justifican ampliamente la inversión. Se recomienda complementar el análisis con una valoración cualitativa de estos beneficios.

\section{Análisis de Viabilidad}

La viabilidad del proyecto se evalúa en múltiples dimensiones para justificar el proyecto SIC-MINCH ante la Gerencia de MINCH SRL.

\begin{longtable}{|L{3cm}|C{1.5cm}|L{10cm}|}
\hline
\rowcolor[gray]{0.9}
\textbf{Tipo de viabilidad} & \textbf{Sí/No} & \textbf{Justificación} \\
\hline
Técnica & Sí & Equipo con experiencia en Laravel, Livewire, MySQL. El stack es maduro, bien documentado y con amplia comunidad. El servidor VPS puede alojar la aplicación con los requisitos identificados. \\
\hline
Operativa & Sí & Los 3 usuarios clave (Auxiliar, Contador, Gerente) pueden adaptarse con capacitación. Los procesos actuales son manuales (Excel), por lo que la mejora es evidente y la curva de aprendizaje es baja gracias a la interfaz web intuitiva con TallStackUI. \\
\hline
Económica & Parcial & El VAN a 5 años es negativo considerando solo ahorros directos (-60.461 Bs.). Sin embargo, los beneficios intangibles (información en tiempo real, reducción de errores, cumplimiento tributario) y la extensión del horizonte a 6--7 años hacen viable el proyecto. \\
\hline
Legal & Sí & El sistema cumple con la normativa tributaria boliviana (RC-IVA, IUE, IT) al implementar los cálculos según las tasas oficiales. Los registros electrónicos tienen validez legal. No se almacenan datos sensibles protegidos por leyes especiales. \\
\hline
Cronológica & Sí & 4 meses es un plazo realista para la construcción del proyecto. Las actividades están desglosadas y priorizadas. El equipo puede cumplir con las entregas parciales usando la metodología Scrum con sprints de 2 semanas. \\
\hline
De recursos humanos & Sí & El personal interno de MINCH puede dedicar horas parciales (8--15\%) sin afectar sus funciones principales. Los técnicos externos se dedican full-time al desarrollo. \\
\hline
De riesgos & Media & Los riesgos son controlables con las mitigaciones propuestas. El riesgo más crítico (pérdida de datos) se mitiga con backup automático diario. El aumento de alcance se controla con el comité de cambios semanal. \\
\hline
\end{longtable}

\section{Requerimientos de Hardware y Software Adicionales}

Este apartado es crítico para la viabilidad técnica y presupuestal. Se detallan los requerimientos de infraestructura para desarrollo y producción.

\begin{longtable}{|L{4cm}|L{5cm}|C{1.5cm}|L{2.5cm}|C{1.5cm}|C{1.5cm}|}
\hline
\rowcolor[gray]{0.9}
\textbf{Componente} & \textbf{Especificación mínima} & \textbf{Cant.} & \textbf{Uso} & \textbf{Costo unit. (Bs.)} & \textbf{Costo total (Bs.)} \\
\hline
Servidor VPS (producción) & 4 vCPU / 8 GB RAM / 100 GB SSD / 2 TB transferencia & 1 & Base de datos, backend, frontend, reportes & 350/mes $\times$ 12 = 4.200/año & 4.200 \\
\hline
Servidor de respaldos (NAS) & Synology DS220+ con 2 discos de 4 TB & 1 & Copias de seguridad diarias & 3.500 & 3.500 \\
\hline
PC para desarrollo (Técnicos) & Intel i7 / RAM: 32 GB / SSD 512 GB & 2 & Desarrollo, pruebas locales & 7.000 & 14.000 \\
\hline
PC para Auxiliar (pruebas) & Intel i5 / RAM: 16 GB / SSD 256 GB & 1 & Pruebas de usuario (UAT) & 5.000 & 5.000 \\
\hline
PC para Gerente/Contador & Equipos existentes de MINCH & --- & Acceso web al sistema & 0 & 0 \\
\hline
UPS (respaldo eléctrico) & 1500VA / 900W & 2 & Servidor y NAS & 1.200 & 2.400 \\
\hline
Switch de red + cableado & Gigabit 8 puertos + cables Cat6 & 1 & Conectividad local & 700 & 700 \\
\hline
\rowcolor[gray]{0.9}
\textbf{Total Hardware} & & & & & \textbf{29.800} \\
\hline
\end{longtable}

\begin{longtable}{|L{4cm}|L{3cm}|L{2cm}|C{1.5cm}|C{2cm}|C{2cm}|}
\hline
\rowcolor[gray]{0.9}
\textbf{Software} & \textbf{Versión / Tipo} & \textbf{Licenciamiento} & \textbf{Cant.} & \textbf{Costo unit. (Bs.)} & \textbf{Costo total (Bs.)} \\
\hline
Sistema operativo servidor & Ubuntu Server 24.04 LTS & Libre & 1 & 0 & 0 \\
\hline
Base de datos & MySQL 8.0 & Libre & 1 & 0 & 0 \\
\hline
Backend runtime & PHP 8.3 + Laravel 13 & Libre & 1 & 0 & 0 \\
\hline
Servidor web & Nginx & Libre & 1 & 0 & 0 \\
\hline
Respaldo automatizado & Duplicati o Veeam Free & Libre & 1 & 0 & 0 \\
\hline
IDE para desarrollo & Visual Studio Code & Libre & 2 & 0 & 0 \\
\hline
Control de versiones & Git + GitHub privado & Gratuito & 1 & 0 & 0 \\
\hline
VPN para acceso remoto & WireGuard & Libre & 1 & 0 & 0 \\
\hline
\rowcolor[gray]{0.9}
\textbf{Total Software} & & & & & \textbf{0} \\
\hline
\end{longtable}

\begin{longtable}{|L{4cm}|L{5cm}|L{6cm}|}
\hline
\rowcolor[gray]{0.9}
\textbf{Requerimiento} & \textbf{Especificación mínima} & \textbf{Justificación} \\
\hline
Ancho de banda de internet & 10 Mbps simétricos & Acceso remoto del Gerente, reportes por email \\
\hline
Red interna & 100 Mbps garantizados & Comunicación entre servidor y clientes \\
\hline
IP fija o DDNS & IP pública fija o servicio DDNS & Acceso externo para Gerente y Contador remoto \\
\hline
Firewall & Router con NAT y reglas básicas & Seguridad perimetral \\
\hline
\end{longtable}

\textbf{Costo estimado de implementación de red (si no existe):} 2.500 Bs.

\begin{longtable}{|L{4cm}|C{2cm}|C{2cm}|C{2cm}|C{2.5cm}|}
\hline
\rowcolor[gray]{0.9}
\textbf{Equipo} & \textbf{Consumo estimado (W)} & \textbf{Horas/día} & \textbf{kWh/mes} & \textbf{Costo mensual (Bs.)} \\
\hline
Servidor VPS (en la nube) & --- & 24 & --- & Incluido en hosting \\
\hline
NAS (respaldo local) & 50 W & 24 & 36 kWh & 56,52 \\
\hline
PC desarrollo (2) & 300 W c/u & 8 & 144 kWh & 226,08 \\
\hline
PC Auxiliar & 150 W & 6 & 27 kWh & 42,39 \\
\hline
Switch + router & 30 W & 24 & 21,6 kWh & 33,91 \\
\hline
\rowcolor[gray]{0.9}
\textbf{Total mensual} & & & \textbf{228,6 kWh} & \textbf{358,90} \\
\hline
\end{longtable}

*Costo por kWh = 1.57 Bs. estimado*

\begin{longtable}{|L{8cm}|C{4cm}|}
\hline
\rowcolor[gray]{0.9}
\textbf{Concepto} & \textbf{Costo (Bs.)} \\
\hline
Hardware (servidores, PCs, UPS, switch) & 29.800 \\
\hline
Software base (open source) & 0 \\
\hline
Red y conectividad (implementación) & 2.500 \\
\hline
Consumo eléctrico anual (proyección) & 4.307 \\
\hline
\rowcolor[gray]{0.9}
\textbf{Total inversión en infraestructura} & \textbf{36.607} \\
\hline
\end{longtable}

\end{document}
